%	현대물리실험 보고서
%	실험4
%	202100973 이승엽

%----------------------------------------------------------------------------------------
%	PACKAGES AND DOCUMENT CONFIGURATIONS
%----------------------------------------------------------------------------------------

\documentclass[a4paper, 10pt, nanum]{CSUniSchoolLabReport}
% use UTF8 encoding
\usepackage[utf8]{inputenc}
% use KoTeX package for Korean
\usepackage{kotex}
% use ams' math font
\usepackage{amsmath}
\usepackage{hyperref}
% use table
\usepackage{graphicx}
\usepackage{multirow}
\usepackage{indentfirst}
\usepackage{setspace}
\usepackage{enumitem}
\usepackage{wrapfig}
\usepackage{epstopdf}
% use captions
\usepackage{caption}
\usepackage{subcaption}
\usepackage{blindtext}
% use layout frame
\usepackage{showframe}

% only for show page layout
\renewcommand\ShowFrameLinethickness{0.25pt}
\renewcommand*\ShowFrameColor{\color{red}}

\setlength{\parindent}{0.2in} % 들여쓰기 길이 설정
\setlength{\parskip}{3mm} % 문단 간의 간격 조절
\setstretch{1.5} % 줄간격
% \graphicpath{{Figures/}} % fig 경로 설정

\addbibresource{reference.bib} % Bibliography file (located in the same folder as the template)

%----------------------------------------------------------------------------------------
%	REPORT INFORMATION
%----------------------------------------------------------------------------------------

\title{현대물리실험 실험4 보고서 \\ Scanning Tunneling Microscope} % Report title

\author{\textsc{Department of Physics} 202100973 이승엽}

\date{\today}

%----------------------------------------------------------------------------------------

\begin{document}

\maketitle % Insert the title, author and date using the information specified above

\begin{center}
	\begin{tabular}{l r}
		Date Performed: & May 16, 2023 \\ % Date the experiment was performed
		& May 23, 2023 \\
		Partners: & 202100969 이규리 \\ % Partner names
		& 202100989 한누리 \\
		Instructor: & Professor 이기주 \\ % Instructor/supervisor
		Typesetting: & LaTeX \\
		Datafitting: & Python \\
	\end{tabular}
\end{center}

%----------------------------------------------------------------------------------------
%	ABSTRACT
%----------------------------------------------------------------------------------------

\maketitle
% \begin{abstract}
% 	This report ...
% \end{abstract}

%----------------------------------------------------------------------------------------
%	INTRODUCTION
%----------------------------------------------------------------------------------------

\section{Introduction}

  To

%----------------------------------------------------------------------------------------
%	THEORY
%----------------------------------------------------------------------------------------

\section{Theory}

	% \begin{figure}[htb!]
	% 	\centering
	% 	\includegraphics[width=5cm]{fig1.png}
	% 	\caption{Top : Simplified term diagram for neon.\\ Bottom : the electron current flowing to the collector as a function of the acceleration voltage in the Franck-Hertz	experiment with neon.}
	% 	\label{fig:1}
	% \end{figure}

  J

%----------------------------------------------------------------------------------------
%	EXPERIMENTAL METHOD
%----------------------------------------------------------------------------------------

\section{Experimental Method}

	% \begin{figure}[htb!]
	% 	\centering
	% 	\includegraphics[width=5cm]{fig2.png}
	% 	\caption{Schematic diagram of the Franck-Hertz tube, Ne.}
	% 	\label{fig:2}
	% \end{figure}
d

%----------------------------------------------------------------------------------------
%	RESULT AND DISCUSSION
%----------------------------------------------------------------------------------------

\section{Result and Discussion}

% \subsection{Franck-Hertz curve:}
% 	\begin{figure}[htb!]
% 		\centering
% 		\includegraphics[width=8cm]{fig6.pdf}
% 		\caption{Data fitting for Franck-Hertz curve.}
% 		\label{fig:6}
% 	\end{figure}

	Fig. 6은 0 V부터 80 V까지 $U_2$를 조절하여 전하의 흐름 데이터를 fitting한 것이다. 그래프를 분석하여 Tab. 1으로 정리한 것이다.

	\begin{table}[htb!]
		\label{tab:1}
		\centering
		\caption{전하의 흐름이 상승(peak)하는 부분이다.}
		\begin{tabular}{c|c}
				\noalign{\smallskip}\noalign{\smallskip}\hline\hline
				$n$ & $U_2$ (V) \\
			\hline
				1 & 19 \\
				2 & 36 \\
				3 & 55 \\
			\hline
			\hline
		\end{tabular}
	\end{table}

	이론적인 부분과 비교했을 때, 실 데이터와 약 30\% 차이가 나지만 그 간격과 그래프 모양을 보면 차이가 많지 않음을 볼 수 있다.

% \subsubsection{Light emission:}
% 	\begin{figure}[htb!]
% 		\centering
% 		\includegraphics[width=5cm]{fig7.1.png}
% 		\caption{Light emission}
% 		\label{fig:7}
% 	\end{figure}

	\begin{table}[htb!]
		\label{tab:2}
		\centering
		\caption{전하의 흐름이 상승(peak)하는 부분이다.}
		\resizebox{\textwidth}{!}{
		\begin{tabular*}{\textwidth}{@{\extracolsep{\fill}}ccc|cc}
				\noalign{\smallskip}\noalign{\smallskip}\hline\hline
				& $n$ & & $U_2$ (V) \\
			\hline
			& 1 & & 29.9 \\
			& 2 & & 47.4 \\
			& 3 & & 65.7 \\
			\hline
			\hline
		\end{tabular*}
		}
	\end{table}

	이론적인 부분과 비교했을 때, 오차가 1\% 내이기에 이론값과 차이가 없음을 알 수 있다.

%----------------------------------------------------------------------------------------
%	CONCLUSIONS
%----------------------------------------------------------------------------------------

\section{Conclusions}

	Franck-Hertz experiment with neon에서 네온 원자에 속박된 전자가 양자화되어 있음을 확인할 수 있었다. 특히 Fig. 6을 보면 특정 peak에서 전류가 상승하는 것을 볼 수 있고, 이는 전하가 이동함을 뜻한다. 즉, 그 순간에 속박에서 풀려난다는 것을 알 수 있다. 또한, Fig. 7을 통해 orbit을 건널 때마다 에너지, 즉 빛이 방출됨을 볼 수 있었다.

%----------------------------------------------------------------------------------------
%	FITTING CODE FOR PYTHON
%----------------------------------------------------------------------------------------

\section{Fitting code for Python}

% \begin{figure*}[htb!]
% 	\centering
% 	\includegraphics[width=8cm]{code1.png}
% 	\label{fig:code}
% \end{figure*}

%----------------------------------------------------------------------------------------
%	BIBLIOGRAPHY
%----------------------------------------------------------------------------------------

% \printbibliography % Output the bibliography

%----------------------------------------------------------------------------------------

\end{document}